\begin{figure}[tb]
\centering
\begin{tikzpicture}[scale=1.05,annotation text]
\coordinate (O) at (0,0);
\coordinate (A) at (3.0,1.5);
\coordinate (B) at (1.65,.35);
\draw[coordinate axis,reference quantity,strong line] (O)--(A) node[above]{\(\boldsymbol\psi_s\)};
\draw[coordinate axis,guide quantity,strong line] (O)--(B) node[below]{\(L_q\vect i_s\)};
\draw[coordinate axis,derived quantity,strong line] (B)--(A) node[midway,right]{\(\boldsymbol\psi_a\)};
\node[draw,rounded corners,derived light region,align=center] at (6.4,.8)
 {\(\boldsymbol\psi_s=L_q\vect i_s+\boldsymbol\psi_a\)\\
 remove the isotropic response\\
 retain torque-producing flux};
\end{tikzpicture}
\caption{Active-flux vector subtraction. The remainder is the flux of the
equivalent nonsalient torque model.}
\label{fig:active-flux-subtraction}
\end{figure}
